% ============================================================================
% Formato.tex: configuración común de los informes del Laboratorio de
% Circuitos Eléctricos II (LIE-311). Se incluye desde Documento.tex.
% La portada se omite intencionalmente y se agrega por separado.
% ============================================================================

% Codificación, idioma y tipografía. inputenc mantiene compatibilidad con
% pdfLaTeX; en LuaLaTeX/XeLaTeX puede omitirse.
\usepackage[utf8]{inputenc}
\usepackage[T1]{fontenc}
\usepackage[spanish,es-nodecimaldot]{babel}
\usepackage{lmodern}

% Matemáticas, unidades, tablas e imágenes.
\usepackage{amsmath,amssymb}
\usepackage{siunitx}                 % Ej.: \SI{10}{\volt}
\usepackage{graphicx}
\usepackage{booktabs}
\usepackage{array}
\usepackage{enumitem}
\usepackage{float}
\usepackage{caption}
\usepackage{xcolor}

% Diseño de página, separación de párrafos y enlaces.
\usepackage{geometry}
\usepackage{fancyhdr}
\usepackage{parskip}
\usepackage{csquotes}
\usepackage[hidelinks]{hyperref}

% El estilo numérico produce citas compactas, apropiadas para IEEE. Para APA,
% cambie style=numeric-comp por style=apa (y agregue el paquete correspondiente
% si su instalación de biblatex no lo incluye).
\usepackage[
    backend=biber,
    style=numeric-comp,
    sorting=none
]{biblatex}

% La comprobación permite compilar la plantilla antes de crear references.bib.
% Al entregar, ese archivo debe contener el manual y todas las fuentes usadas.
\IfFileExists{references.bib}{\addbibresource{references.bib}}{}

\geometry{a4paper, margin=2.5cm, headheight=15pt}
\setlength{\parindent}{0pt}
\setlength{\parskip}{6pt}
\setlist{itemsep=2pt, topsep=3pt}
\captionsetup{font=small, labelfont=bf}
\sisetup{locale=DE, per-mode=symbol, detect-all}

% Datos editables. Se redefinen en Documento.tex con \renewcommand. Se declaran
% antes del encabezado porque este último los utiliza.
\newcommand{\numeroPractica}{N}
\newcommand{\nombrepractica}{Título de la práctica}

% Encabezado y pie. Los valores de la práctica se definen en Documento.tex.
\pagestyle{fancy}
\fancyhf{}
\lhead{LIE-311 -- Laboratorio de Circuitos Eléctricos II}
\rhead{Práctica \numeroPractica}
\cfoot{\thepage}

% Inserta una figura centrada y numerada.
% Uso: \imagen[etiqueta]{archivo}{pie de figura}
% Las capturas virtuales deben mostrar fecha y hora; el archivo puede ser PNG,
% JPG o PDF. Una etiqueta permite referenciar la figura con \ref{fig:montaje}.
\newcommand{\imagen}[3][]{%
    \begin{figure}[H]
        \centering
        \includegraphics[width=0.88\textwidth]{#2}
        \caption{#3}
        \if\relax\detokenize{#1}\relax\else\label{#1}\fi
    \end{figure}%
}

% Entorno reutilizable para objetivos y conclusiones numerados.
\newenvironment{itemsinforme}
    {\begin{enumerate}[label=\arabic*., leftmargin=*]}
    {\end{enumerate}}

% Marcador visible para localizar contenido pendiente. Debe reemplazarse o
% eliminarse antes de entregar el informe final.
\newcommand{\rellenar}[1]{\textcolor{red}{\textbf{[Completar: #1]}}}

\hypersetup{
    pdftitle={Informe de laboratorio -- \nombrepractica},
    pdfauthor={Nombre del estudiante},
    pdfsubject={Laboratorio de Circuitos Eléctricos II -- LIE-311}
}
